% RESEARCH_DRAFT_SHA256: 69f87b3890383d60636d3fded5f84233db55b66ea85838247ff30f0813e94cd7
% Layout and source-reference wrappers expanded for the web renderer.
\documentclass[11pt]{article}
\usepackage[margin=25mm]{geometry}
\usepackage{amsthm}
\usepackage{fontspec,unicode-math}
\setmainfont{Latin Modern Roman}
\setmathfont{Latin Modern Math}
\usepackage{amsmath,amssymb,mathtools,xurl,hyperref}
\hypersetup{colorlinks=true,urlcolor=blue,linkcolor=blue}
\setcounter{secnumdepth}{0}
\setlength{\parindent}{0pt}
\setlength{\parskip}{5pt}
\setlength{\emergencystretch}{3em}
\newcommand{\sref}[2]{\hyperlink{source:#1:#2}{[S#2]}}
\newcommand{\cataloguescope}[1]{\paragraph{Source record.}#1}
\newcommand{\floor}[1]{\left\lfloor#1\right\rfloor}
\newtheorem{proposition}{Proposition}
\newtheorem{lemma}[proposition]{Lemma}
\begin{document}
\paragraph{Research manuscript.} OpenAI \texttt{gpt-6-astra}, reasoning \texttt{ultra}; 11 October 2026.
\paragraph{Preserved source body.} The complete catalogue mathematical text and sources follow with layout and source links expanded, including source qualifications and its historical preparation notes. New research begins at the section headed Research attempt.

\section{Convolutional Codes for Retrieval, Sensing, and Rank Metrics}\label{statement}

\subsection{Definitions and mathematical statement}

\paragraph{Shared notation.}Standard number systems are $\mathbb N,\mathbb Z,\mathbb Q,\mathbb R,\mathbb C$. Specialized notation and index ranges follow the cited source and the definitions below.


A convolutional encoder maps a symbol stream to code symbols using a fixed finite window of past inputs. Sliding-window recovery concerns erasures or errors within consecutive windows. Private information retrieval seeks a requested database item while hiding its index from the specified servers; compressed sensing seeks sparse vectors from linear measurements. For a vector over $\mathbb F_{q^m}$, rank weight is the rank over $\mathbb F_q$ of its coordinate-expansion matrix, and rank distance is the rank weight of a difference.
\paragraph{Record scope and qualifications.}
The catalogue title is retained as metadata. Named certificates, contracts, bounds or frameworks in that title are not asserted as results of the original source extract.
The supplied question proposes three directions rather than one parameterized construction. Server assumptions, sensing guarantees, and rank-metric target parameters must be specified in a proposed answer.

Source-record limitation: The code families, sliding-window setup, and required retrieval/sensing/rank-metric guarantees are not identified.
\paragraph{Statement recorded in the supplied source.}
Can convolutional codes exploit sliding-window properties for private information retrieval or compressed sensing, and can they attain strong rank-metric properties? \href{https://doi.org/10.4171/owr/2019/13}{[S2]}
\paragraph{Source status.}
The supplied catalogue labels this source record open; that metadata was checked against the supplied snapshot on 2026-10-11. The label is preserved for provenance and is not an independent verification of present open status or mathematical correctness. The catalogue literature assessment is dated 2026-08-22.
\subsection{Short English statement}
Determine how finite-memory codes can support retrieval, sparse recovery, or correction of rank errors.
\subsection{Sources}
\begin{itemize}
\item[S1] Ulam AI and the UnsolvedMath Contributors, UnsolvedMath: A Curated Collection of Open Mathematics Problems, catalogue entry. Catalogue curation and metadata: CC BY 4.0. The source extract is retained; definitions, English summaries and qualifications are editorial additions.
\url{https://huggingface.co/datasets/ulamai/UnsolvedMath}.
\item[S2] Contemporary Coding Theory (2020), Oberwolfach Reports, DOI 10.4171/owr/2019/13, pp. 773--840
\url{https://doi.org/10.4171/owr/2019/13}.
\item[S3] Convolutional rank-metric codes
\url{https://arxiv.org/abs/1404.7251}.
\item Source-specific Oberwolfach Reports licences, including CC BY-SA 4.0 where applicable, are retained; see THIRD\_PARTY\_NOTICES.md.
\end{itemize}
\paragraph{Preparation and review.}
The original source and UnsolvedMath contributors retain attribution for the source extract. Codex (GPT-6; exact serving revision not exposed) assisted with definitions, English summaries, scope notes and formatting. No new human mathematical authorship or independently verified solution is asserted. Catalogue identity checked against the complete supplied registry and current repository. Distinct upstream records, including repeated questions and broad or damaged extracts, are retained for complete coverage of the source CS category. A repeated record is not a new independent conjecture. Complete mathematical-equivalence and current-literature review has not been performed. Preparation and final review follow contribution prompts C and D, with source gaps disclosed. Structural and rendering checks do not establish a new theorem or open status.
\label{end:statement}

\section{Research attempt}
\subsection{Literature, source context and scope}
The full original catalogue body and source qualifications are preserved above. The relevant primary workshop passages were read in the cached publisher report. Historical claims and subsequent attributed results are distinguished from the mathematical work below. A restricted example, standard lemma or known theorem is not counted as new progress on the entire original target. Missing assumptions and repeated or rhetorical extracts are not silently turned into different solved problems.

\subsection{A sliding-window inverse needs a boundary anchor}
The source lists retrieval, compressed sensing and rank-metric coding as separate directions. We test a basic convolutional recovery step that might be used in a sliding-window construction. In a real-valued model, suppose
\[
y_t=u_t+u_{t-1},\qquad 1\le t\le L,
\]
where $L\ge1$. This is a scalar finite-memory filter, not by itself a private-retrieval or compressed-sensing scheme.
\begin{proposition}
If $u_0$ is known, the window is uniquely recoverable by
\[
u_t=(-1)^tu_0+\sum_{j=1}^t(-1)^{t-j}y_j.
\]
If each observed $y_j$ has additive error of magnitude at most $\eta$, the worst possible error in $u_t$ is exactly $t\eta$. If $u_0$ is unknown, the window measurements alone do not determine the input sequence.
\end{proposition}
\begin{proof}
The recursion $u_t=y_t-u_{t-1}$ proves the inverse formula by induction. With an exact initial anchor, the error equals $\sum_{j=1}^t(-1)^{t-j}e_j$, whose magnitude is at most $t\eta$. Taking $e_j=(-1)^j\eta$ makes all summands have the same sign for every fixed $t$ and attains that bound. If the boundary anchor is unknown, add $\delta u_t=(-1)^t a$ for any nonzero real $a$ and every $0\le t\le L$. Then $\delta u_t+\delta u_{t-1}=0$, so all measurements stay unchanged while the input changes.
\end{proof}
\subsection{Verification and remaining gaps}
Exact rational inversions were checked for several windows, along with alternating perturbations attaining the $t\eta$ error and the unknown-boundary null direction. Thus an overlapping-window recovery argument must account both for carried state and for accumulated noise; local invertibility alone supplies no uniform stability bound independent of window length. There are $L$ measurements for $L$ unknown inputs once the anchor is fixed, so this calculation gives no compression advantage. It also supplies no privacy guarantee or rank-metric distance construction. Those separate source goals remain unaddressed beyond this concrete recovery prerequisite, and no full-target progress is claimed.

\section{References and source attribution}
Primary workshop report: \url{https://doi.org/10.4171/owr/2019/13}. Relevant passages read: 2019/13 p. 794, separate convolutional PIR, compressed-sensing and rank-metric directions.

All retrievals were checked on 11 October 2026. Original question attribution remains in the preserved Sources section. External theorems are attributed and not claimed as new discoveries or fully verified merely by citation.


\paragraph{Integration review.} Codex (GPT-6; exact serving revision not exposed) reviewed the written arguments and expanded layout and source-reference wrappers. Repository metadata and the production web renderer were checked using Python 3.13 and JavaScript with MathJax 3.2.2. This model-assisted review is not independent human proof verification; no PDF compilation or formal proof checking is asserted.

\section{Completion Estimate}
\noindent\textbf{COMPLETION: 0\%}
No substantiated new progress toward the entire original target is claimed. The worked results and failed route are delimited above; a known result, illustrative model, or auxiliary obstruction is not counted as original-target completion.

\section{Final status and attribution}
\textbf{Attempt status: unresolved.} This describes this research attempt, not a new ruling on the source catalogue's status or on external literature claims.
The mathematics and this self-audit were generated by OpenAI \texttt{gpt-6-astra}, reasoning \texttt{ultra}, on 11 October 2026. Codex (GPT-6; exact serving revision not exposed) orchestrated the batch, reviewed the written mathematical arguments, and contributed formatting/repository review. This additional model review is not independent human verification. No human mathematical authorship, independent human proof checking, or formal verification is claimed. The historical preparation paragraphs in the reproduced catalogue body concern that body only.

Contribution prompt A was applied to the complete definition and supplied source record, followed by the adversarial checks of prompt B. Local reproduction files are \texttt{batch022/checks.py}, \texttt{computations.json}, and the per-ID audit JSON. They were not public links at preparation time. Numerical checks are finite and do not prove unrestricted results. No LaTeX engine was found locally; compilation was not performed by this research agent.

\end{document}
